\dropcap{T}{he stone} had been warm when the midwife put it against him and was warm no longer.

It lay under his belly in the fold of the cloth, a small river stone the colour of skimmed milk, giving up its heat with the slow reluctance of a thing that has been asked and has agreed. It had been warmed at a neighbour's fire --- at the kiln-house down the yard, where somebody kept banked coals through the cold --- and carried here in a rag, still hot from a hearth that was not this hearth. That was a favour. It was not entered anywhere. \person{Daniel Marchetti} had known households where a kindness goes in warm and comes out written in a hand you cannot read, and men who warm a stone for a cot and never ask who paid for the fire. He could not have said which kind this was. He could only feel the stone going cold against him, a little at a time, the way he had once felt a patient's core temperature fall under his palm.

The child was \person{Corin}, two hours old, and called nothing yet that he could answer to. He lay on a fold of grey wool in the front room of \place{House Silt}, one of two clay-daubed rooms that opened off the \place{Clay Yard} in \place{Low Kiln}, the poor quarter of \place{Sill}, in \place{Grath}. The floor was trodden ash, packed by generations of feet until it had the dull spring of old bread, and along the skirting board ran a grey ridge the width of a thumb: \lore{the Ash}, sifted out of the walls and the pots and the lives that had been kept here, left because nobody sweeps it before the Five. The Five were three days off. Outside, the \place{Ilm} lay under its skin of ice, and the cold came through the door that would not shut and went around the room looking for a place to sit.

It smelled of wet wool and tallow and clay dust, and under those the sour mineral edge of the Ash, which is the smell of a room that has been warm a long time and has kept things it should have let go. A woman sang somewhere in the house, a low song with the words flattened out by work, and the tune went under the skin and stayed. He understood the words. He could not make one.

\person{Alda Reed} worked with steady hands. She had the stoop of a woman who had spent her life leaning over other people's births, and the forearms of one who had lifted other people's children; her nails were scrubbed to the quick and still had clay under them, a thin brown line that no soap had ever got out. She addressed the stone first, turning it in the rag with two fingers the way a cook tests a pot, then the slip pot on the hearth, then the child, then the mother. She did not address the father, who was not in the room and would not be addressed. This was not cruelty. It was the correct number of words for a room where somebody was still in pain.

The slate lay within reach of her right hand: a flat grey slab, a dusting of chalk, a line of writing in a strange neat hand that Dan did not look at twice because he had other things to look at and no way yet to look at anything. Her hand went to the slate before it went to the child. He registered the movement and filed it where he filed everything he could not use yet, under the heading of things that would matter later.

“Warm the stone first,” she said, and her voice was low and even, the voice of a woman who has said the same sentence four hundred times and never once said it carelessly. “A cold stone is a stone, not a kindness.” She said it to the stone, or to the room, or to \person{Dessa Silt}, who was in no condition to argue. Then something about the slip, about not setting it near the fire, about paying for it in the thaw --- practical instruction addressed to the pot before it was addressed to a person, because the pot would answer and the person would not.

Dessa lay with her face turned toward the wall and her breath coming in a rasp that Dan knew before he knew that he knew it: expiration long, inspiration short, a small catch at the end that was not relief but the body trying again. Her right hand went to her left forearm, to a pale scar the width of a sieve-wire, and rubbed it in a slow circular press while she counted under her breath.

“Nineteen,” she said. “Eighteen. Nineteen.”

She had walked those nineteen steps every day for six years, from the slip pit to the drying rack, and counted them then for the length of the working day, for the bits at the end of it, for the number of times her feet would go before she could stop. Now she was counting them for a different reason and Dan could not know what it was. He could only hear the count and know, with the part of him that had once stood under the lights of a trauma bay and sorted twenty people in ninety seconds, that she was not counting steps at all. She was holding something out of the room by counting it in.

His eyes --- his new eyes, wet and unfocused, not yet good at holding a face --- went to the angle of her wrist instead of to her face. The scar sat where sieve-wire had cut it, and the way she pressed it told him she was deciding something she did not want to decide. He read the pulse at the wrist before he understood that he was reading it: quick, thready, a little irregular at the end of expiration. He read the colour under the nails, grey-blue at the beds, and the fine tremor in the hand, which was not from cold. He counted the exits without deciding to --- the door to the yard, the window with its oiled paper, the gap under the lintel where light came in thin --- and then counted them again to check the count. This was what his mind did now. It had nothing to do, so it took the room apart.

He watched Alda's hands because they were the steadiest thing in the room. They did not hurry and they did not hesitate; they went from the stone to the slip to the cloth to the mother in an order that had been decided years before Dan was born and would not be changed now. He had worked with surgeons whose hands could not be hurried and with surgeons whose hands could not be trusted, and he could have told you which was which within a minute of watching them tie a knot. Alda tied nothing, but the order of her reaching was the same kind of signature. When she touched the little tin at her bodice --- once, briefly, before she went back to the child --- his eye caught the gesture and held it. He did not know what it was. He knew that it was a thing done every time and never for the first time.

The stone was giving up its heat at a rate he could almost measure. Not fast --- a stone gives heat the way a bank gives change, reluctantly, a little at a time --- but steadily, and the rag between them was damp now where his body had sweated against it, and damp cloth draws heat faster than dry. He had watched patients cool in a white room far away and had known, without a thermometer, the exact point at which the surface warms and the core falls: \lore{the Afterdrop}, the worsening that arrives wearing the face of improvement. He knew it the way a priest knows sin. There was no rewarming here, no blanket but the scrap of wool, no patient but himself, and still he knew the shape of it --- knew that the warmth going out of the stone was going into a body that could not hold it, and that in a few minutes the body would be colder than it had been before the stone was put there.

He knew the Afterdrop when he saw it. He wanted the breast. Both were true and the second was louder.

The body that wanted was not the body that knew. The knowing body was a large soft man with forearms like a butcher's and a left knee that clicked in the cold, a man who had reduced a thousand shoulders in the dark and could find the thin place under skin with two fingers and a name for it. The wanting body was two hours old, had never held anything up, and was cold and hungry and furious about it. It drooled. It could not hold up its head. When it tried to turn over it stayed where it was, and the grey wool that Alda had tucked around him had come loose at one corner and would not go back. He had set a femur in a ditch in the rain and he could not keep a scrap of wool on his own shoulder.

The joke was there, dry and small, and it cost him something. He did not laugh, because laughing is a thing the body does with air and the body had other plans for its air. It wanted. That was all it had ever wanted, and it wanted now with a whole man's urgency in a vessel the size of a loaf.

He could not speak. He could not reach a cup. A trauma surgeon who could intubate in blood and in the dark, who could cut a way into a destroyed throat with a knife and a tube, lay in a fold of cloth and could not lift his own hand to his own mouth. There was a cup somewhere on the hearth --- he had seen Alda set it down --- and the distance to it was two paces, which is a long way when you have no hands for the journey. He could not have said the word for cup in Grathic. He could not have said it in English either, because the body had not yet learned to push air into words. He had a whole ledger in his head and no way to open it.

It was the first small grief of the new life, and it was not for the dead, or for the daughter in a country that had no name here, or for the call he had not made. It was for a pair of hands that could do nothing. He had spent twelve years being useful in the only way he knew how, which was to decide quickly and carry the decision, and here was a room that needed nothing decided and a body that would not do what it was told. The competence was intact. There was simply nothing behind it. It sat in him like an instrument in a dead man's hand.

Outside, something moved on the frozen river: not the river, but the ice under it, or the wind over it, a small sound like a hand rubbing dry paper. Dan listened to it because listening was what he had left. He counted the song against Dessa's breathing and got a ratio that meant nothing and kept it anyway. He had been counting things for as long as he could remember, pulses and ribs and steps from the bed to the door in rooms he would never sleep in again, and he had never once been told that it was a gift. It was a habit. It was the way he had learned to stand in a room where he could not help.

A woman came in from the yard and stood just inside the door with the cold still on her sleeves. \person{Mira Silt}. Eel slime had dried on the wool in grey patches, and she smelled of river silt and salted perch and the cold off the water. There was something in her pocket that ticked when she moved her thumb, small and hard; she flicked it once against a nail and put it away. She did not sit. There was a stool by the wall and a clear place on the table where a slate might be laid, and she went near neither.

“Hold him?” she said, and it was not clear whether the question was for anyone. “You want me to. You don't. I can see which one's winning.” She looked at the child, or at the place where the child was, and did not reach for him. Then she said nothing further, and stayed standing.

Alda finished with the slip and wiped her hands on a rag that had been a shirt once. She lifted the child --- not roughly, but as a woman lifts a sack of meal, with the whole weight taken in one movement --- and set him back against the folded cloth, and in doing so turned him so that the stone slid a little and lay closer to the hollow of his belly. The grey wool went over him again. This time it stayed for a moment. Then an arm came out of it, small and brown and determined, and the wool was off at the corner again.

“Keep him covered,” Dessa said, and her voice went flat and finished: nothing more to be said on it. She reached over and pressed the wool down at the corner with two fingers, pressing as if to seal something, and let her hand lie there a moment longer than the cloth required. The pale scar crossed the back of her forearm in the tallow light. Then she took her hand away and went back to counting, and what she was holding out of the room went back in.

Mira took her hand off the door frame and went out the way she had come, leaving the smell of silt and salt behind her and the small ticking stopped. She had not held him. She had not sat down. Dan watched her go with the part of him that watched people leave rooms and noted, without knowing why, that she had kept her back to the slate the whole time.

Alda went back to the slate and stood over it with the chalk in her fingers, and the room settled around the standing of her body the way water settles around a stone dropped in it. The banked fire on the hearth shifted and settled, and a small flame came up under the ash and went down again. She did not look at it. She wrote something on the slate, blew the dust off it, and set the slate down within reach of her hand again.

Dan listened to the room the way he had once listened to a corridor at four in the morning: for the note under the noise. The fire had a small irregular tick as the clay of the flue cooled. Dessa's breath had its rasp and its catch. Alda's skirt dragged on the ash floor when she moved. Under all of it, under the song that had stopped and the count that went on, there was the stone, and it had a sound too if you put your ear to it --- not a sound in the air but a sound in the body, the small settling of warmth leaving a thing that will not be warmed again. He had heard that sound in patients. He had never heard it in himself.

Under the child the stone had reached the point where it no longer gave heat but only failed to take it. The warmth that had been in it was in him now, or had gone into the room, and the room was warm around him with its fire and its wet wool and its breathing bodies, and nobody in it was cold. Nobody in it felt the stone go. Alda's hand was on the slate. Dessa was counting under her breath. Mira was out in the yard with the cold still on her sleeves. And under the grey wool, in the hollow of a belly that did not yet know it was hollow, the stone went on giving up the last of what it had --- quietly, without a sound, and with nobody the wiser.
